\usepackage{graphicx}
\usepackage{geometry}
% \usepackage{natbib}
\geometry{
    top=30mm,
    bottom=30mm,
    left=30mm,
    right=30mm
}
\usepackage{setspace}

% Layout robustness: reduce overfull boxes from long technical prose.
% This does not change content; it allows TeX a bit more flexibility in line breaking.
\setlength{\emergencystretch}{2em}

\usepackage{titlesec}

% Use XeLaTeX or LuaLaTeX (fontspec). Do not use pdfLaTeX with this config.
\usepackage{fontspec}

\usepackage[hidelinks]{hyperref}
\usepackage{xurl}
\usepackage{booktabs}
\usepackage{tabularx}
\usepackage{amsmath, amssymb}
\usepackage{amsfonts}
\usepackage{cleveref}

\usepackage{tikz}
\usetikzlibrary{shapes,arrows,positioning,calc,patterns,decorations.markings}
\usetikzlibrary{shapes,arrows,positioning,calc,patterns,decorations.pathreplacing}

\usetikzlibrary{shapes,arrows,positioning,calc,arrows.meta,backgrounds}
\usetikzlibrary{shadows}
\usetikzlibrary{decorations.pathmorphing}
% \usepackage{epigraph}


% Times New Roman is often missing on Linux; use a Times-like fallback.
\IfFontExistsTF{Times New Roman}{%
  \setmainfont{Times New Roman}%
}{%
  \IfFontExistsTF{TeX Gyre Termes}{%
    \setmainfont{TeX Gyre Termes}[
      BoldFont = * Bold,
      ItalicFont = * Italic,
      BoldItalicFont = * Bold Italic,
    ]%
  }{%
    \setmainfont{Liberation Serif}%
  }%
}

% Chapter formatting
\titleformat{\chapter}[block]  % "block" puts label + title together
  {\centering\bfseries\fontsize{15pt}{18pt}\selectfont} % formatting: centered, bold, 15pt font
  {Chapter \thechapter} % label text
  {1em} % spacing between label and title
  {} % code before title text

% Section formatting
\titleformat{\section}
  {\bfseries\fontsize{14pt}{16pt}\selectfont} % bold + 14pt font, with line spacing
  {\thesection} % section number
  {1em} % spacing between number and title text
  {} % code before the title text

\titleformat{\subsection}[block]
  {\bfseries\fontsize{13pt}{15pt}\selectfont} % bold + 13pt font
  {\thesubsection} % number
  {1em}            % spacing between number and text
  {}             % actual title text

% ===== THAI UNICODE SUPPORT =====
% Don't use polyglossia - it can interfere with font selection
% Just define Thai font directly
\newfontfamily\thaifont{THSarabunNew.ttf}[
    Path = ./font/,
    BoldFont = THSarabunNew Bold.ttf,
    ItalicFont = THSarabunNew Italic.ttf,
    BoldItalicFont = THSarabunNew BoldItalic.ttf,
    Script = Thai
]

% Simple Thai command that only changes font for Thai text
\newcommand{\thai}[1]{{\thaifont #1}}
% ================================